\section{Introduction}
\label{sec:intro}

Autism Spectrum Disorder (ASD) is a neurodevelopmental condition affecting
roughly 1--2\,\% of the global population, characterised by wide variation
in social communication, learning, and behaviour. Its prevalence and
lifelong impact have made it a central focus of genetic research.

\subsection{The Genetic Basis of ASD}

ASD is highly heritable: estimates suggest genetics account for around
80\,\% of the risk. However, no single gene is responsible for the
disease itself. Instead, ASD arises from a combination of hundreds of
small-effect common variants and rarer, higher-impact mutations, many of
which influence brain development and synaptic function. This dual
genetic architecture, common and rare variants, is typically studied in
isolation, making a unified understanding difficult.

Genome-Wide Association Studies (GWAS) scan the genomes of thousands of
individuals to identify variants more common in people with ASD. While
they have flagged over 100 genomic regions of interest, the findings
have critical limitations:

\begin{itemize}
  \item Each associated variant explains only a tiny fraction of risk on
        its own.
  \item Polygenic Risk Scores (PRS) estimate an individual's genetic
        predisposition to a specific disease or trait by aggregating the
        effects of thousands or millions of small genetic variants
        across their entire genome. They provide a personalised,
        numerical risk score, often expressed as a percentile.
        Polygenic Risk Scores built from known GWAS variants explain
        less than 5\,\% of why autism runs in families.
  \item Standard GWAS tests each variant in isolation, missing the
        gene--gene interactions that likely play a meaningful role.
  \item Statistical signals are difficult to translate into biological
        insight: knowing a genomic region is associated with ASD does
        not readily reveal which gene is involved or how.
\end{itemize}

Despite advances, the known genetic variants from GWAS still leave the
vast majority of ASD's heritability unexplained. Bridging the gap
between statistical association signals and genuine biological
understanding requires new analytical strategies, ones that can handle
high-dimensional data, respect the interconnected nature of gene
function, and leverage the growing wealth of biological knowledge
databases. Machine learning, applied thoughtfully to GWAS data,
represents the most promising current direction for meeting that
challenge.
